\chapter{Analisis temporal y costes de desarrollo}\label{anatemporal}

\section{Planificacion temporal}

La planificacion de BargAIn se definio con una dedicacion objetivo de 300 horas,
distribuidas en fases funcionales y tecnicas. El proyecto se ejecuto de forma iterativa,
manteniendo trazabilidad en \texttt{TASKS.md} y en el arbol \texttt{.planning/}.

\begin{table}[h]
\centering
\caption{Resumen por fases del proyecto}
\label{tab:fases-temporales}
\begin{tabular}{|l|c|c|c|}
\hline
\textbf{Fase} & \textbf{Semanas} & \textbf{Horas est.} & \textbf{Estado} \\
\hline
F1 Analisis y diseno & S1--S3 & 45 & Completada \\
F2 Infraestructura base & S3--S5 & 30 & Completada \\
F3 Desarrollo core backend & S5--S12 & 95 & Completada \\
F4 Desarrollo frontend & S8--S16 & 70 & Completada \\
F5 IA, optimizador y scraping & S10--S17 & 35 & Completada \\
F6 Pruebas, deploy y cierre & S17--S20 & 25 & Completada \\
F7 Cierre final defensa & Extra & 25 & Completada \\
\hline
\textbf{Total} & & \textbf{300} & \textbf{Cerrado} \\
\hline
\end{tabular}
\end{table}

\section{Metodologia de trabajo}

Se aplico una metodologia incremental con ciclos cortos por fase:

\begin{itemize}
	\item definicion de alcance y criterios de aceptacion,
	\item implementacion por modulos,
	\item validacion automatizada (lint + tests),
	\item verificacion funcional y documental,
	\item cierre de fase con artefactos de evidencia.
\end{itemize}

Este enfoque permitio absorber cambios sin comprometer estabilidad del sistema.

\section{Estimacion de costes}

La estimacion economica se realizo mediante coste-hora por rol junior,
alineada con referencias del informe de perfiles TIC de la Junta de Andalucia (2018).

\begin{table}[h]
\centering
\caption{Coste estimado del desarrollo}
\label{tab:coste-estimado}
\begin{tabular}{|l|c|c|}
\hline
\textbf{Bloque} & \textbf{Horas} & \textbf{Coste (EUR)} \\
\hline
Analisis y diseno & 45 & 1.490,40 \\
Infraestructura & 30 & 861,60 \\
Backend + frontend & 165 & 4.738,80 \\
IA, optimizador y scraping & 35 & 1.005,20 \\
Pruebas y despliegue & 25 & 594,25 \\
\hline
\textbf{Total} & \textbf{300} & \textbf{8.690,25} \\
\hline
\end{tabular}
\end{table}

Adicionalmente, el coste operativo en staging se mantuvo bajo durante el TFG
gracias al uso de planes gratuitos o de bajo consumo (Render, Sentry, GitHub Actions).

\section{Riesgos y mitigaciones}

Los riesgos operativos principales fueron:

\begin{itemize}
	\item \textbf{Cambios en fuentes de scraping:} mitigado con spiders modulares y
		fallback de crowdsourcing.
	\item \textbf{Dependencia de APIs externas:} mitigado con cache y degradacion controlada.
	\item \textbf{Carga de trabajo en fases finales:} mitigado con cierre incremental de
		entregables y automatizacion de validacion.
\end{itemize}

\section{Resultado de planificacion}

La planificacion permitio completar el alcance funcional priorizado del TFG,
incluyendo arquitectura, implementacion, pruebas y despliegue de staging,
manteniendo coherencia entre roadmap, requisitos y evidencias de verificacion.


